\documentclass[10pt]{article}
\usepackage[top=0.67in,bottom=0.72in,left=0.83in,right=0.83in]{geometry}
\usepackage[cmex10]{amsmath}
\usepackage{amsthm}
\usepackage[hidelinks]{hyperref}
\DeclareFontShape{OMX}{cmex}{m}{n}{<->cmex10}{}
\DeclareMathSizes{10}{10}{7}{7}
\DeclareMathSizes{9}{10}{7}{7}
\linespread{0.955}
\newtheorem{theorem}{Theorem}
\title{Incomparability constraints for a 55-member Frankl counterexample}
\author{Achita Anantachina}
\date{26 November 2026}
\begin{document}
\maketitle
\begin{abstract}
We consider the proposed first remaining size for a minimum
counterexample to the union-closed sets conjecture. In each of the
four 55-member cases, any two maximal points in the reduced domination
order have at least four common incomparable points. The domination
order must have at least 20 incomparable pairs.
\end{abstract}

\section{The four cases}
Frankl's conjecture asks whether every finite union-closed family with
a nonempty member has a point in at least half the members. Suppose a
smallest counterexample has 55 members and choose one with smallest
support among those. The 12-element theorem gives support size at
least 13 \cite{vuckovic}. Write \(m\) for that size. We use the maximal
avoiding set argument of Hu and its formulation in \cite{bhattacharjee}.

The chosen family is separating. If two points appeared in exactly
the same members, projecting away one would preserve distinctness,
union closure, and every remaining frequency. A minimal member can
be removed while preserving union closure. Thus at least one point
\(t\) has frequency 27, since otherwise that removal would give a
smaller counterexample. Let \(A_i\) be the union of the members avoiding
\(i\), and let \(U\) be the support. Define \(x\preceq y\) when every
member containing \(x\) also contains \(y\). Separation makes this a
partial order. The \(A_i\) are distinct and \(t\in A_i\) whenever
\(i\ne t\), since otherwise \(t\preceq i\) would violate separation
or the frequency cap.

The 28 members avoiding \(t\) form a smaller union-closed family, so
some point occurs in at least 14 of them. Choose such a point \(p\)
maximal under domination among those with that property. If
\(i\notin\{t,p\}\) and \(p\notin A_i\), then \(p\preceq i\), so \(i\)
also has the property. Therefore \(p\) belongs to \(U\) and all
\(A_i\) for \(i\in W=U\setminus\{t,p\}\). In particular, \(p\)
occurs in at least \(14+m-1\) members and at most 27, so \(m\le14\).

Split members into four blocks by membership of \(t,p\). Let \(R,S,T\)
be the sets of their complements within \(W\) for blocks 10, 01, and
00. Put \(D(i)=\{x\in W:x\preceq i\}\). The block 11 complements
include the empty mask and all \(D(i)\), with at most one extra mask
\(E\). If \(b\) is the number of block 11 members and \(a=|S|\), then
\(b\ge m-1\), \(a\ge14\), and \(a+b\le27\). The other sizes are
\(|R|=27-b\) and \(|T|=28-a\). This gives
\[
\begin{array}{c|rrr|c}
m & |R| & |S| & |T| & \text{extra mask}\\ \hline
14 &14&14&14&\text{no}\\
13 &15&14&14&\text{no}\\
13 &15&15&13&\text{no}\\
13 &14&14&14&\text{yes}
\end{array}
\]

All masks are ideals of the domination order on \(W\). A singleton
member gives an injection from sets avoiding its point into sets
containing it, by taking unions. For a two-point member the same
injection takes sets containing neither point to sets containing both,
so the two frequencies sum to at least the family size. Thus no
member has size one or two. If \(w=|W|\), all masks in \(R,S\) have
size at most \(w-2\). Masks in \(T\) have size at most \(w-3\) or
are all of \(W\), and every block 11 mask is proper. Each point of
\(W\) belongs to at least 28 complement masks, counted over all four blocks.

\section{The incomparability bound}
\begin{theorem}
For any of the four cases, any two maximal points of \(W\) have at
least four common incomparable points, and \(W\) has at least
20 incomparable pairs.
\end{theorem}
\begin{proof}
There are at least two maxima, since a greatest point would have
principal ideal \(W\), giving the forbidden member \(\{t,p\}\).
For distinct maxima \(x,y\), let \(C=D(x)\cup D(y)\) and
\(h=w-|C|\). A mask containing both contains \(C\), and \(h\) is the
number of common incomparable points. The full member \(U\) and
the \(w-2\) members corresponding to principal ideals other than
\(D(x),D(y)\) contain both points. The two classes of members
avoiding \(x\) and \(y\) therefore have union of size at most
\(55-(w-1)\). Their sizes are each at least 28, so at least \(w\)
complement masks contain both points.

If \(h\le3\), at most four such masks occur in each of \(R,S\), at
most two in \(T\), and at most one in block 11, namely \(E\) if present.
Here the four possibilities for \(R,S\) are \(C\) and its extensions
by one point. The two possibilities for \(T\) are \(C\) and \(W\).
This total is below \(w\), except possibly when \(w=11\), \(h=3\), and
\(E\) is present and all the bounds are attained. Then both \(R\)
and \(S\) contain all four permitted masks, and their cross
intersections include distinct masks containing \(x,y\). Such an
intersection must be a block 11 mask by union closure, and the
only possibility is \(E\). This is a contradiction. Thus \(h\ge4\).
In particular, there cannot be only two maxima since every point
lies below some maximal point.

Let \(k\) be the number of maxima and \(N_i\) the nonmaximal points
incomparable with maximum \(i\). For \(k=3,4,5\), each pair of such sets
has intersection of size at least \(6-k\). A point of \(N_i\)
is comparable to at least one maximum other than \(i\), so it belongs
to at most \(k-2\) of the other sets. Hence
\[
|N_i|\ge\left\lceil\frac{(k-1)(6-k)}{k-2}\right\rceil.
\]
Counting these pairs and the maximal antichain gives at least
21, 18, and 20 incomparable pairs for \(k=3,4,5\), respectively.
For \(k\ge7\), the maximal antichain alone suffices.

For \(k=4\), suppose there are at most 19 incomparable pairs.
There are at least 12 between maxima and nonmaxima, so at most one
lies entirely among nonmaxima. If the nonmaxima form a chain, their
incomparability sets \(N_i\) are nonempty terminal segments of that
chain. The last point of the chain would lie in every \(N_i\).
Every nonmaximal point lies below some maximal point, so this is
impossible. If there is one nonmaximal
incomparable pair, then \(|N_i|=3\) for each maximum. Each \(N_i\)
has at least two points in common with each other \(N_j\), so all its
points must occur in exactly three of the four sets. The 12
memberships give four points, each below a different sole maximum.
They must be pairwise incomparable, since comparability would force
one to lie below a second maximum. This gives six nonmaximal pairs.

For \(k=6\), at most 19 incomparable pairs allow at most four outside
the maximal antichain. If each nonmaximal point is incomparable with
at most three maxima, every ideal containing four maxima contains
all nonmaxima. Thus a mask in \(R\) or \(S\) includes at most four maxima
and a proper mask in \(T\) at most three. The full mask in \(T\)
contributes at most six, the principal masks together contribute
six, and \(E\) contributes at most five. The total over the six maxima
is at most the following quantity, where \(\epsilon=1\) if \(E\) exists
and zero otherwise.
\[
4(|R|+|S|)+3(|T|-1)+12+5\epsilon.
\]
The four bounds are 163, 167, 168, and 168. The only other case has
one nonmaximal point incomparable with four maxima, and every other
nonmaximal point lies below all six maxima. The bounds increase by
at most one, since a proper mask in \(T\) with four maxima must simply
omit that point and take precisely the four maxima incomparable
with it. In either case each of \(R,S\) falls at least one below its
individual cap. If all masks had four maxima, a trace on the maxima
would determine at most two masks. There are at least 14 masks in
each family, so two have different traces and their intersection is
a member with at most three maxima. Subtracting two from even
the increased bound gives less than the required \(6\cdot28=168\).
\end{proof}

\section{Computational search}
The companion encoder uses Boolean variables for the domination
relation and the masks in each block. It imposes cardinality,
separation-based frequency bounds, and the incomparability restrictions
above. For pairs of blocks, a union of members corresponds to an
intersection of masks in the block determined by their label union.
A full encoding adds all such conditions. The lazy variant adds the
condition for a pair of row positions when the current model violates
it. A separate integer-mask checker verifies every union and frequency
of any extracted family.

The bounded runs reported in the companion data do not settle any
of the four full cases. The general conjecture remains open.
The computation and exposition used AI assistance.

\begin{thebibliography}{9}
\bibitem{vuckovic} B. Vu\v{c}kovi\'c and M. \v{Z}ivkovi\'c.
\href{https://ipsitransactions.org/journals/papers/tir/2017jan/fullPaper.pdf}
{The 12-element case of Frankl's conjecture} (2017).
\bibitem{bhattacharjee} D. Bhattacharjee and U. Bhattacharya.
\href{https://zenodo.org/records/23645724}
{Attaining the entropy limit for union-closed families} (2026).
\end{thebibliography}
\end{document}
